\chapter{Diskussion und kritische Reflexion}
\label{ch:diskussion}

\section{Einordnung der Ergebnisse}
Die Kurve in Abbildung~\ref{fig:modell} bildet den erwarteten Verlauf
des theoretischen Modells ab. Da keine Messdaten vorliegen, kann weder
die Gültigkeit des Modells experimentell überprüft noch eine
Übereinstimmung mit realen Beobachtungen behauptet werden.

\section{Grenzen, Fehlerquellen und Alternativerklärungen}
Für eine echte Messreihe wären beispielsweise eine ungenaue Bestimmung
der Pendellänge, der Startwinkel und die Zeitmessung zu berücksichtigen.
Ergebnisse müssen mit ihrer Unsicherheit und ihrem Gültigkeitsbereich
bewertet werden, nicht nur mit einer optisch passenden Kurve.

\section{Eigenanteil und Zusammenarbeit}
% Besonders wichtig bei Arbeiten, die am SFN betreut werden:
% Wer hat das Thema entwickelt, Geräte gestellt, gemessen,
% beraten, programmiert, ausgewertet oder Text korrigiert?
In einer tatsächlichen BLL wird hier transparent dokumentiert,
welche Aufgaben selbstständig ausgeführt wurden und welche Unterstützung
beispielsweise durch Lehrkräfte, SFN-Betreuende oder andere Personen
erfolgte. Das ersetzt keine vollständige Quellen- und
Hilfsmitteldokumentation.
